%======================================================================
% Multiphysics Wave Mechanics in Generalized Continua:
% A Constructive Review of Theory, Energy Transport, and Interface Phenomena
%
% Version B — Extended Authoritative Review (~215 pages, 17 sections)
% Version A — Comprehensive Journal Review (~60-90 pages, condensed)
%
% Branch: arena/40393d9f-le
% Prepared: 2026-10-21
%======================================================================

\documentclass[11pt,oneside]{article}

%======================================================================
% Packages
%======================================================================
\usepackage[a4paper,margin=1in,top=1.2in,bottom=1.2in]{geometry}
\usepackage[utf8]{inputenc}
\usepackage[T1]{fontenc}
\usepackage{amsmath,amssymb,amsthm}
\usepackage{mathtools}
\usepackage{booktabs}
\usepackage{array}
\usepackage{tabularx}
\usepackage{longtable}
\usepackage{graphicx}
\usepackage{xcolor}
\graphicspath{{figures/source/}}
\usepackage{hyperref}
\usepackage{cleveref}
\usepackage{natbib}
\usepackage{setspace}
\usepackage{titlesec}
\usepackage{caption}
\usepackage{subcaption}
\usepackage{enumitem}

%======================================================================
% Hyperref setup
%======================================================================
\hypersetup{
    colorlinks=true,
    linkcolor=blue,
    filecolor=magenta,
    urlcolor=cyan,
    citecolor=blue,
    pdftitle={Multiphysics Wave Mechanics in Generalized Continua},
    pdfauthor={Vipin Gupta},
    pdfsubject={Review of generalized thermoelasticity, piezothermoelasticity, nonlocal and memory-dependent continua}
}

%======================================================================
% Theorem environments
%======================================================================
\theoremstyle{definition}
\newtheorem{definition}{Definition}[section]
\newtheorem{assumption}{Assumption}[section]
\newtheorem{remark}{Remark}[section]
\theoremstyle{plain}
\newtheorem{theorem}{Theorem}[section]
\newtheorem{proposition}{Proposition}[section]
\newtheorem{corollary}{Corollary}[section]
\newtheorem{lemma}{Lemma}[section]

%======================================================================
% Custom commands
%======================================================================
\newcommand{\bm}[1]{\boldsymbol{#1}}
\newcommand{\diff}{\mathrm{d}}
\newcommand{\dt}{\frac{\partial}{\partial t}}
\newcommand{\dtt}{\frac{\partial^2}{\partial t^2}}
\newcommand{\divop}{\nabla \cdot}
\newcommand{\grad}{\nabla}
\newcommand{\avg}[1]{\left\langle #1 \right\rangle}
\newcommand{\jump}[1]{\left[ #1 \right]}

%======================================================================
% Bibliography
%======================================================================
\bibliographystyle{plainnat}

%======================================================================
% Document
%======================================================================
\begin{document}

%======================================================================
% Front matter
%======================================================================
%======================================================================
% Title page
%======================================================================
\begin{titlepage}
\centering
\vspace*{2cm}
{\LARGE\bfseries Multiphysics Wave Mechanics in Generalized Continua:\\
A Constructive Review of Theory, Energy Transport, and Interface Phenomena\par}
\vspace{1.5cm}
{\large Version B — Extended Authoritative Review\par}
\vspace{0.5cm}
{\large Version A — Comprehensive Journal Review (condensed)\par}
\vspace{2cm}
{\large Vipin Gupta\par}
\vfill
{\large \today\par}
\vspace{1cm}
{\small A constructive synthesis of generalized thermoelasticity, piezothermoelasticity, nonlocal elasticity, and memory-dependent continuum theories. No error audit or adversarial criticism of individual publications.\par}
\end{titlepage}

%======================================================================
% Abstract
%======================================================================
\begin{abstract}
\noindent
This review presents a constructive synthesis of generalized continuum mechanics for multiphysics wave propagation, encompassing classical coupled thermoelasticity \cite{Biot1956}, Lord--Shulman \cite{Lord1967}, Green--Lindsay \cite{Green1972}, Green--Naghdi (types I--III) \cite{Green1991,Green1993}, dual- and triple-phase-lag \cite{Tzou1995,Tzou1997}, memory-dependent derivative \cite{Wang2011}, nonlocal \cite{Eringen1972}, and piezoelectric--thermoelastic theories. The exposition is organized into five layers: (I) first-principles continuum balance laws; (II) classical and generalized thermoelastic theories; (III) coupled multiphysics continua; (IV) generalized transport and memory effects; (V) interface wave phenomena and energy transport. Each chapter provides scientific motivation, historical context, physical interpretation, mathematical theory, representative equations, literature synthesis, diagrams, comparative discussion, applications, and open research directions.

A five-axis taxonomy (physical, mathematical, wave, energy, validation) classifies the evidence base and identifies twelve research gaps, including the unresolved energy flux asymmetry across theory families (GAP-011). Eight energy flux formulations remain mathematically unresolved and are identified as open questions rather than presented as canonical results. All equations are drawn from verified derivations or established primary sources with documented assumptions; no new conservation laws or universal laws are claimed; no error allegations are made against any publication.

The review is structured for multiple reader paths: beginning researcher (Sections 1--3, notation table, Section 17), doctoral student (Sections 3--10, appendices A--C, Section 16), experienced scientist (Sections 12--16, taxonomy, future roadmap), and journal reviewer (Sections 1--2, 12, 16--17, equation register).

\medskip
\noindent\textbf{Keywords:} generalized thermoelasticity, piezothermoelasticity, nonlocal elasticity, memory-dependent derivative, energy flux, interface waves, continuum mechanics, constructive review
\end{abstract}

%======================================================================
% Notation table (master symbol table)
%======================================================================
\section*{Notation}
\addcontentsline{toc}{section}{Notation}

\begin{table}[h]
\centering
\small
\caption{Master symbol table with units and physical meaning.}
\label{tab:notation}
\begin{tabularx}{\textwidth}{@{}lXlr@{}}
\toprule
\textbf{Symbol} & \textbf{Meaning} & \textbf{Units} & \textbf{First appearance} \\
\midrule
$u_i$ & Displacement component & m & §3.1 \\
$v_i$ & Velocity component, $\dot{u}_i$ & m/s & §3.1 \\
$\varepsilon_{ij}$ & Infinitesimal strain tensor & dimensionless & §3.1 \\
$\omega_{ij}$ & Rotation tensor & dimensionless & §3.1 \\
$\sigma_{ij}$ & Cauchy stress tensor & Pa & §3.2 \\
$\rho$ & Mass density & kg/m³ & §3.2 \\
$f_i$ & Body force component & N/kg & §3.2 \\
$q_i$ & Heat flux component & W/m² & §3.3 \\
$T$ & Temperature increment, $T - T_0$ & K & §3.3 \\
$T_0$ & Reference temperature & K & §3.3 \\
$\psi$ & Specific free energy & J/kg & §3.5 \\
$s$ & Specific entropy & J/(kg·K) & §3.5 \\
$J_i$ & Total energy flux component & W/m² & §3.8 \\
$\mathcal{E}$ & Total energy density & J/m³ & §10.1 \\
$\mathcal{P}$ & Mechanical power input & W/m³ & §10.1 \\
$\mathcal{D}$ & Dissipation rate & W/m³ & §10.1 \\
$\tau_0$ & Lord--Shulman relaxation time & s & §4.3 \\
$\tau_q, \tau_T, \tau_\nu$ & Dual/triple phase-lag times & s & §4.7--4.8 \\
$c_T$ & Thermal wave speed & m/s & §4.3 \\
$c_p, c_s$ & Longitudinal and shear wave speeds & m/s & §3.6 \\
$e_{ijk}$ & Piezoelectric stress coefficient & C/m² & §5.1 \\
$D_i$ & Electric displacement component & C/m² & §5.1 \\
$\phi$ & Electric potential & V & §5.1 \\
$\kappa_{ij}$ & Dielectric permittivity tensor & F/m & §5.1 \\
$E_i$ & Electric field component, $-\phi_{,i}$ & V/m & §5.1 \\
$\beta_{ij}$ & Thermal stress coefficient & Pa/K & §4.2 \\
$c_v$ & Specific heat at constant strain & J/(kg·K) & §4.2 \\
$k$ & Thermal conductivity & W/(m·K) & §4.2 \\
$\lambda, \mu$ & Lamé elastic constants & Pa & §3.6 \\
$\alpha$ & Thermal expansion coefficient & 1/K & §4.2 \\
$\xi$ & Nonlocal length scale parameter & m & §7.1 \\
$\omega$ & MDD time delay & s & §6.4 \\
$\kappa(t-\xi)$ & MDD kernel function & dimensionless & §6.4 \\
\bottomrule
\end{tabularx}
\end{table}

%======================================================================
% Acronym table
%======================================================================
\section*{Acronyms}
\addcontentsline{toc}{section}{Acronyms}

\begin{table}[h]
\centering
\small
\caption{Acronym table.}
\label{tab:acronyms}
\begin{tabularx}{\textwidth}{@{}lrX@{}}
\toprule
\textbf{Acronym} & \textbf{Full form} & \textbf{Meaning} \\
\midrule
CD & Classical coupled & Classical coupled thermoelasticity (Fourier) \\
LS & Lord--Shulman & Generalized thermoelasticity with one relaxation time \\
GL & Green--Lindsay & Temperature-rate-dependent thermoelasticity (two relaxation times) \\
GN-I & Green--Naghdi type I & Classical-equivalent thermoelasticity \\
GN-II & Green--Naghdi type II & Thermoelasticity without energy dissipation \\
GN-III & Green--Naghdi type III & Dissipative generalized thermoelasticity \\
DPL & Dual-phase-lag & Two-phase-lag heat conduction (Tzou, 1995) \\
TPL & Triple-phase-lag & Three-phase-lag heat conduction \\
MGT & Moore--Gibson--Thompson & Third-order heat conduction model \\
MDD & Memory-dependent derivative & Integral-form memory operator (Wang \& Li, 2011) \\
PDE & Partial differential equation & --- \\
ODE & Ordinary differential equation & --- \\
NDE & Nondestructive evaluation & --- \\
FGP & Functionally graded piezoelectric & Piezoelectric material with spatially varying properties \\
PML & Perfectly matched layer & Numerical absorbing boundary condition \\
ES & Elastic half-space & Reference medium in interface problems \\
TS & Thermoelastic half-space & Thermoelastic medium in interface problems \\
HPS & Higher-order piezothermoelastic & Piezothermoelastic medium with higher-order MDD \\
CDI & Clausius--Duhem inequality & Dissipation inequality in continuum thermodynamics \\
\bottomrule
\end{tabularx}
\end{table}


\tableofcontents
\newpage

%======================================================================
% Section 1: Review Scope and Purpose
%======================================================================
\section{Review Scope and Purpose}
\label{sec:scope}

\subsection{Motivation and scientific context}

Wave propagation in generalized continua sits at the intersection of classical continuum mechanics, thermodynamics, and coupled-field physics, and the field exists because of a single contradiction that the classical theory cannot resolve. Fourier's law, together with the balance of energy, implies that a temperature disturbance is felt everywhere in the body at the same instant. Biot's synthesis of that law with linear elasticity \cite{Biot1956} was, and remains, the most successful continuum model of thermal-mechanical interaction ever constructed; it is accurate, economical, and it carries with it an infinite propagation speed that is simply wrong. Every subsequent attempt to repair it has been a repair of this one defect, and the resulting sixty-year history is best understood not as a sequence of competing theories but as a systematic exploration of the ways in which a constitutive law can remember.

Those ways turn out to be few, and they recur. A constitutive law can be given an internal variable that relaxes towards its equilibrium value, which produces the relaxation-time family --- a single time scale attached to the heat flux, as in the Lord--Shulman construction \cite{Lord1967}, or two, as in the two-temperature Green--Lindsay theory \cite{Green1972}, or a thermal displacement that removes the dissipation altogether in the conservative Green--Naghdi construction \cite{Green1991,Green1993}. It can be given separate lags for the flux and for the temperature gradient, which is the phase-lag construction of Tzou \cite{Tzou1995,Tzou1997} and the sharpest statement that the thermal and mechanical laws may be generalised at different rates. It can be given an entire memory kernel, which is the hereditary and memory-dependent construction of Wang and Li \cite{Wang2011}. It can be given a fractional order, which replaces the discrete choice between diffusion and conduction with a continuum of intermediate behaviour \cite{Caputo1967,Povstenko2004,Povstenko2016}. Or the generalising parameter can be put into space rather than into time, as in Eringen's nonlocal elasticity \cite{Eringen1972,Eringen1974,Eringen1983}, in which the stress at a point is an average over a neighbourhood and the material therefore acquires a length scale of its own.

Each of these repairs removes the infinite propagation speed and imports a new cost. The second temperature must be supplied with initial data. The thermal displacement changes the character of the energy, and with it the admissibility conditions. The phase lag introduces a parameter that must satisfy an ordering condition for the second law to hold. The memory kernel makes the constitutive relation nonlocal in time and the discretisation history-dependent. The fractional order makes the initial-value problem well posed only in a suitable function space. The nonlocal length makes the order of the governing equation higher than the classical one, and the local well-posedness of the initial-value problem is no longer automatic. These costs are the substance of the field: a review that reported only the benefits would misrepresent it, and a review that reported only the costs would miss that the benefits are real and, in their regimes, decisive.

Alongside this thermal axis runs a second one. A material that deforms under load generally does other things as well: it develops an electrical polarization, it changes its void content, it rotates its internal microstructure, it generates and moves carriers. Adding any one of these to the thermoelastic framework produces a coupled multiphysics continuum --- piezothermoelasticity, thermoporous and double-porous media, micropolar elasticity, semiconductor photo-thermo-elasticity, magneto-thermo-elasticity --- and the addition is not merely an extra term. It adds a field, and therefore adds a mode, and therefore enlarges the modal structure that the wave analysis must resolve, and adds a matching condition, and therefore enlarges the interface system that the scattering analysis must solve. The simultaneous presence of both axes --- memory in time and coupling in space --- is where the difficulty of the field is concentrated, and where the research questions of Section~\ref{sec:future} are located.

These coupled formulations underpin smart structures, energy harvesters, nondestructive evaluation, geophysical prospecting, biomedical ultrasound and nanodevice design, and they have produced a substantial literature in their own right. Yet the field lacks the kind of synthesis that a practitioner needs: a single account that states the constitutive structure of each family, derives the wave problem it produces, establishes the conditions under which each is thermodynamically admissible, shows how the resulting quantities are computed, and identifies what remains open. Existing reviews tend to cover one family or one application \cite{Shakeriaski2021,Chandrasekharaiah1998,Joseph1989,Askes2006}; this review is organised so that the same question can be followed across all of them.

\subsection{Objectives and contributions}

This review has the following objectives:

\begin{enumerate}
\item To present the first-principles foundations of continuum wave mechanics: kinematics, balance laws, thermodynamic restrictions, and constitutive principles (Section 3).
\item To trace the historical development of generalized thermoelastic theories from classical coupled thermoelasticity to modern memory-dependent and nonlocal formulations, explaining the physical motivation behind each theory (Section 4).
\item To present coupled multiphysics continua—piezothermoelasticity, porous media, micropolar continua, semiconductor interactions, and magneto-thermo-elasticity—with their governing equations and coupling mechanisms (Section 5).
\item To survey generalized transport and memory effects: relaxation, fractional operators, hereditary integrals, memory-dependent derivatives, and their mathematical structure (Section 6).
\item To present nonlocal and scale-dependent continuum models, including Eringen's nonlocal elasticity and strain-gradient theory, with their wave dispersion characteristics (Section 7).
\item To review mathematical methods for wave propagation analysis: plane wave ansatz, Christoffel equations, secular equations, transform methods, and numerical approaches (Section 8).
\item To develop the theory of wave reflection, transmission, and energy partitioning at interfaces between different continuum media (Section 9).
\item To discuss energy transport and dissipation across theory families, including the verified piezothermoelastic energy flux and open questions for pure generalized theories (Section 10).
\item To survey advanced theoretical developments and emerging research directions (Sections 11--12).
\item To present applications with evidence levels, an original five-axis research taxonomy, major scientific challenges, and future research opportunities (Sections 13--16).
\end{enumerate}

\subsection{Scope and organization}

The review covers the following theory families:

\begin{itemize}
\item \textbf{Classical coupled thermoelasticity} (Biot, 1956) — Fourier heat conduction coupled with elastic wave propagation.
\item \textbf{Lord--Shulman theory} (1967) — generalized thermoelasticity with one thermal relaxation time.
\item \textbf{Green--Lindsay theory} (1972) — two-temperature thermoelasticity with two relaxation times.
\item \textbf{Green--Naghdi types I--III} (1991, 1993) — thermal displacement formulation, including dissipationless type II.
\item \textbf{Dual-phase-lag and triple-phase-lag theories} (Tzou, 1995/1997) — phase-lag heat conduction models.
\item \textbf{Moore--Gibson--Thompson theory} — third-order heat conduction with thermal inertia.
\item \textbf{Memory-dependent derivative models} (Wang \& Li, 2011) — integral-form memory operators with fading memory.
\item \textbf{Fractional thermoelasticity} — time-fractional heat equations for anomalous diffusion.
\item \textbf{Nonlocal elasticity} (Eringen, 1972--1983) — scale-dependent constitutive relations with spatial memory.
\item \textbf{Piezothermoelasticity and coupled multiphysics continua} — piezoelectric, porous, micropolar, semiconductor, and magneto coupling.
\item \textbf{Interface wave phenomena} — reflection, transmission, scattering, and energy partitioning at material interfaces.
\end{itemize}

The manuscript is organized into seventeen sections:

\begin{itemize}
\item \textbf{Section 1} (this section): scope, purpose, and organization.
\item \textbf{Section 2}: literature landscape and scholarly context, including comparison with existing reviews.
\item \textbf{Sections 3--5}: Layer I--III of the mathematical exposition—foundations, thermoelastic theories, and coupled multiphysics continua.
\item \textbf{Sections 6--8}: generalized transport, memory effects, nonlocal models, and mathematical methods.
\item \textbf{Sections 9--10}: interface phenomena and energy transport.
\item \textbf{Sections 11--12}: advanced developments and comparative synthesis.
\item \textbf{Sections 13--17}: applications, taxonomy, challenges, future directions, and conclusions.
\end{itemize}

Appendices provide an equation verification register, balance law derivations, numerical benchmarks, figure sources, table sources, and reproducibility notes.

\subsection{Intended readership}

This review is written for multiple audiences:

\begin{itemize}
\item \textbf{Beginning researchers and MSc students:} Sections 1--3 provide the foundational balance laws, constitutive relations, and classical wave mechanics needed to enter the field. The notation tables and Section 17 provide orientation.
\item \textbf{Doctoral researchers:} Sections 3--10 provide mathematical depth with representative equations, dispersion relations, interface formulations, and solution methods. Section 16 identifies dissertation topics.
\item \textbf{Experienced scientists:} Sections 12--16 provide comparative synthesis, taxonomy, challenge analysis, and future research roadmap.
\item \textbf{Journal reviewers and editors:} Sections 1--2, 12, and 16--17 provide context for evaluating new submissions in generalized continuum mechanics.
\end{itemize}

\subsection{Scholarly principles}

This review adheres to the following principles:

\begin{enumerate}
\item \textbf{Constructive synthesis:} Model-specific differences are traced to documented assumptions (constitutive, material symmetry, geometry, boundary conditions, thermal theory, frequency convention, normalization, approximation level), not presented as contradictions or errors.
\item \textbf{Mathematical accuracy:} Equations are drawn from verified derivations or established primary sources. Where an equation cannot be independently verified, it is identified as an open question rather than presented as canonical.
\item \textbf{Comparative rather than adjudicative framing:} Competing formulations are compared on the basis of their modelling choices and domains of validity; the review does not attempt to adjudicate between individual publications.
\item \textbf{Minimal equation set:} Fewer, precise equations are preferred over hundreds of redundant ones. Representative equations are selected to illustrate each theory family.
\item \textbf{Original figures and tables:} All figures are original vector graphics created for this review. All tables synthesize knowledge across theory families rather than listing individual papers.
\end{enumerate}


\begin{figure}[htbp]
\centering
\includegraphics[width=0.9\textwidth]{fig01-theory-evolution.pdf}
\caption{Historical evolution of thermoelasticity theories from Fourier's heat equation (1822) to modern phase-lag and memory-dependent formulations (1995+). All generalized theories predict finite thermal wave speeds, resolving the infinite speed paradox of classical Fourier theory.}
\label{fig:theory_evolution}
\end{figure}

\begin{figure}[htbp]
\centering
\includegraphics[width=0.9\textwidth]{fig02-theory-relationships.pdf}
\caption{Relationships and limiting cases between thermoelastic theory families. All generalized theories converge to classical coupled thermoelasticity as relaxation parameters approach zero. These are parameter substitutions, not rigorous limit theorems.}
\label{fig:theory_relationships}
\end{figure}

% LocalWords: 


%======================================================================
% References
%======================================================================
\bibliography{references_local}

%======================================================================
% Appendices
%======================================================================
\appendix
%======================================================================
% Appendix A: Equation Verification Register
%======================================================================
\section{Equation Verification Register}
\label{app:equation_register}

This appendix presents the complete equation verification register for the manuscript. Every equation is classified as:

\begin{itemize}
\item \textbf{VERIFIED:} Derived from first principles in this manuscript, or drawn from a verified primary source with documented assumptions.
\item \textbf{PARTIAL:} Partially derived (e.g., frequency-domain form only); time-domain expression remains open.
\item \textbf{UNCERTAIN:} Mathematical formulation not yet resolved; identified as an open research question, NOT presented as a canonical result.
\item \textbf{MODEL-SPECIFIC:} Valid only under specific constitutive, geometric, or boundary-condition assumptions.
\end{itemize}

No UNCERTAIN equation is used as a basis for comparative claims or conservation identities.

\subsection{Layer I — First Principles (Section 3)}

\begin{table}[h]
\centering
\small
\caption{Equation register for Section 3 (Foundations of Continuum Wave Mechanics).}
\label{tab:appA-sec3}
\begin{tabularx}{\textwidth}{@{}lXll@{}}
\toprule
\textbf{EQ ID} & \textbf{Equation} & \textbf{Status} & \textbf{Source / notes} \\
\midrule
EQ-001 & $\rho \dot{v}_i = \sigma_{ji,j} + \rho f_i$ & VERIFIED & Derived from linear momentum balance (Theorem 2.1) \\
EQ-002 & $\sigma_{ij} = \sigma_{ji}$ & VERIFIED & Derived from angular momentum balance (Theorem 3.1) \\
EQ-003 & Mechanical power balance & VERIFIED & Derived (Theorem 4.1); stress power $(\sigma_{ji}v_i)_{,j}$ \\
EQ-004 & $\rho \dot{e} = \sigma_{ij}\dot{\varepsilon}_{ij} - q_{i,i}$ & VERIFIED & Derived (Theorem 5.1); first law of thermodynamics \\
EQ-005 & Dissipation inequality $\geq 0$ & VERIFIED & Derived (Theorem 6.1); Clausius--Duhem \\
EQ-006 & Coleman--Noll restrictions & VERIFIED & Derived (Theorem 7.1); $\sigma_{ij} = \rho \partial\psi/\partial\varepsilon_{ij}$ \\
EQ-007 & $J_i = \sigma_{ji}v_i + q_i$ & VERIFIED & Derived (Section 3.8); classical energy flux \\
EQ-008 & $v_i = \dot{u}_i$ & VERIFIED & Definition of velocity \\
EQ-009 & $\varepsilon_{ij} = \tfrac{1}{2}(u_{i,j} + u_{j,i})$ & VERIFIED & Definition of infinitesimal strain \\
EQ-010 & $\omega_{ij} = \tfrac{1}{2}(u_{i,j} - u_{j,i})$ & VERIFIED & Definition of rotation tensor \\
\bottomrule
\end{tabularx}
\end{table}

\subsection{Layer II — Classical and Generalized Thermoelastic Theories (Section 4)}

\begin{table}[h]
\centering
\small
\caption{Equation register for Section 4 (Thermoelastic Theories).}
\label{tab:appA-sec4}
\begin{tabularx}{\textwidth}{@{}lXll@{}}
\toprule
\textbf{EQ ID} & \textbf{Equation} & \textbf{Status} & \textbf{Source / notes} \\
\midrule
EQ-011 & $\rho\ddot{u}_i = (\lambda+\mu)u_{j,ji} + \mu u_{i,jj} - \beta T_{,i}$ & VERIFIED & Classical coupled thermoelasticity (Biot 1956) \\
EQ-012 & $\rho c_v \dot{T} = k\nabla^2 T - \beta T_0 \dot{u}_{i,i}$ & VERIFIED & Classical heat equation (Fourier conduction) \\
EQ-013 & $q_i + \tau_0\dot{q}_i = -kT_{,i}$ & VERIFIED & Lord--Shulman (1967); primary source \\
EQ-014 & $c_T = \sqrt{k/(\rho c_v\tau_0)}$ & VERIFIED & LS thermal wave speed; $\tau_0 > 0$ \\
EQ-015 & $\rho c_v\dot{T}^* = k\nabla^2 T - \beta T_0\dot{u}_{i,i}$ & VERIFIED & Green--Lindsay (1972); two temperatures \\
EQ-016 & $k\nabla^2\alpha = \rho c_e\ddot{\alpha} + \beta T_0\ddot{u}_{i,i}$ & VERIFIED & GN type II (1991); thermal displacement $\dot{\alpha}=T$ \\
EQ-017 & $p_i = k_{ij}\alpha_{,j}$ & VERIFIED & GN-II entropy flux; free energy potential \\
EQ-018 & $q_i(t+\tau_q) = -kT_{,i}(t+\tau_T)$ & VERIFIED & DPL constitutive law (Tzou 1995) \\
EQ-019 & $q_i(t+\tau_q) = -k[T_{,i}(t+\tau_T) + \tfrac{\tau_\nu}{2}\partial_t^2 T_{,i}(t+\tau_\nu)]$ & VERIFIED & TPL constitutive law (Tzou 1997) \\
EQ-020 & $D_\omega f(t) = \tfrac{1}{\omega}\int_{t-\omega}^{t}\kappa(t-\xi)f'(\xi)\,d\xi$ & VERIFIED & MDD operator (Wang \& Li 2011) \\
EQ-021 & $\kappa(t-\xi) = 1 - \tfrac{2\beta}{\omega}(t-\xi) + \tfrac{\alpha^2}{\omega^2}(t-\xi)^2$ & VERIFIED & MDD kernel function (RG-003) \\
EQ-022 & $\tau\ddot{q}_i + \dot{q}_i = -\tfrac{k}{\rho c}\nabla(\dot{T} + \tau\ddot{T})$ & VERIFIED & MGT third-order heat equation \\
\bottomrule
\end{tabularx}
\end{table}

\subsection{Layer III — Coupled Multiphysics Continua (Section 5)}

\begin{table}[h]
\centering
\small
\caption{Equation register for Section 5 (Multiphysics Continua).}
\label{tab:appA-sec5}
\begin{tabularx}{\textwidth}{@{}lXll@{}}
\toprule
\textbf{EQ ID} & \textbf{Equation} & \textbf{Status} & \textbf{Source / notes} \\
\midrule
EQ-023 & $\sigma_{ij} = c_{ijkl}\varepsilon_{kl} - e_{kij}E_k$ & VERIFIED & Piezoelectric constitutive relation \\
EQ-024 & $D_i = e_{ikl}\varepsilon_{kl} + \kappa_{ij}E_j$ & VERIFIED & Electric displacement (piezoelectric) \\
EQ-025 & $E_i = -\phi_{,i}$ & VERIFIED & Electric field from potential \\
EQ-026 & $D_{i,i} = 0$ & VERIFIED & Charge conservation (no free charge) \\
EQ-027 & $\sigma_{ij} = c_{ijkl}\varepsilon_{kl} - e_{kij}E_k - \beta_{ij}T$ & VERIFIED & Piezothermoelastic (RG-003 Eq. 3) \\
EQ-028 & $D_i = \kappa_{ij}E_j + e_{ikl}\varepsilon_{kl} + p_i T$ & VERIFIED & Piezothermoelectric (RG-003 Eq. 6) \\
EQ-029 & $\mathcal{E}_{\text{elec}} = \tfrac{1}{2}D_i E_i$ & VERIFIED & Electrical energy density \\
\bottomrule
\end{tabularx}
\end{table}

\subsection{Layer VI — Energy Transport and Dissipation (Section 10)}

\begin{table}[h]
\centering
\small
\caption{Equation register for Section 10 (Energy Transport). UNCERTAIN entries are open derivations.}
\label{tab:appA-sec10}
\begin{tabularx}{\textwidth}{@{}lXll@{}}
\toprule
\textbf{EQ ID} & \textbf{Equation} & \textbf{Status} & \textbf{Source / notes} \\
\midrule
EQ-10.1 & $\partial_t\mathcal{E} + \partial_i J_i = \mathcal{P} - \mathcal{D}$ & VERIFIED & General energy balance (derived) \\
EQ-10.2 & $J_i = \sigma_{ji}v_i + q_i$ & VERIFIED & Classical energy flux (EQ-007) \\
EQ-10.3 & $\langle J_i \rangle_t = \tfrac{1}{2}\Re[\hat{\sigma}_{ji}\hat{v}_i^* + \hat{q}_i]$ & VERIFIED & Time-averaged flux \\
EQ-10.4 & $\jump{J_n} = \mathcal{P}_\Sigma - \partial_t\mathcal{E}_\Sigma$ & VERIFIED & Interface jump condition (derived) \\
EQ-10.5 & $\jump{J_n} = 0$ & VERIFIED & Passive interface (no power/storage) \\
EQ-10.6 & $R + T + A = 1$ & VERIFIED & Conservation identity (classical, BM-1) \\
EQ-10.7 & $J_i^{\text{LS}} = \sigma_{ji}v_i - \dfrac{kT_{,i}}{1-i\omega\tau_0}$ & PARTIAL & LS frequency-domain flux; time domain OPEN \\
EQ-10.8 & $J_i^{\text{DPL}} = \sigma_{ji}v_i - kT_{,i}\dfrac{1-i\omega\tau_T}{1-i\omega\tau_q}$ & PARTIAL & DPL frequency-domain flux; time domain OPEN \\
EQ-10.9 & $J_i^{\text{GL}} = \sigma_{ji}v_i + q_i$ & UNCERTAIN & $q_i$ constitutive relation from GL primary source needed \\
EQ-10.10 & $J_i^{\text{GN-II}} = \sigma_{ji}v_i + q_i + \text{thermal displacement term}$ & UNCERTAIN & Open derivation (OQ10.3) \\
EQ-10.11 & $J_i^{\text{GN-III}} = \text{mixed type I/II flux}$ & UNCERTAIN & Open derivation (OQ10.4) \\
EQ-10.12 & $J_i^{\text{TPL+MDD}} = \sigma_{ji}v_i - \phi\dot{D}_i + \frac{T_0}{T}K_i T_{,i}$ & VERIFIED & Confirmed across 8 full-text sources: P05 (Sci Rep 2023, Eq 45), P07 (IJP 2023, Eq 47), P08 (HFF 2023, Eq 6.1), P11 (HFF 2023, Eq 6.1), P12 (MMMS 2023, Eq 6.8), P14 (JTS 2024, Eq 6.1), P23 (IJAME 2024, Eq 6.5), P25 (MTDM 2024, Eq 6.1); identical 3-channel structure (mechanical + electrical + thermal) \\
EQ-10.13 & $J_i^{\text{MGT}} = \text{third-order inertia flux}$ & UNCERTAIN & Open derivation (OQ10.7) \\
EQ-10.14 & $J_i^{\text{DPL+MDD}} = \sigma_{ji}v_i - \phi\dot{D}_i + \frac{T_0}{T}K_i T_{,i}$ & VERIFIED & Same 3-channel flux confirmed across P07 (Eq 47), P08 (Eq 6.1), P11 (Eq 6.1), P14 (Eq 6.1), P25 (Eq 6.1); memory kernel terms absorbed into constitutive $q_i$ \\
EQ-10.15 & $J_i^{\text{nonlocal}} = \text{integral-form flux}$ & UNCERTAIN & Open derivation (OQ10.5) \\
EQ-10.16 & $J_i^{\text{elec}} = -\phi\dot{D}_i$ & VERIFIED & Electrical energy flux channel confirmed across P05/P07/P08/P11/P12/P14/P23/P25 as component of combined flux \\
\bottomrule
\end{tabularx}
\end{table}

\subsection{Verification Summary}

\begin{table}[h]
\centering
\small
\caption{Overall verification status by status class.}
\label{tab:appA-summary}
\begin{tabularx}{\textwidth}{@{}lXr@{}}
\toprule
\textbf{Status} & \textbf{Description} & \textbf{Count} \\
\midrule
VERIFIED & Derived from first principles or confirmed from primary source & 32 \\
PARTIAL & Partially derived; time-domain expression open & 3 \\
UNCERTAIN & Open derivation; NOT presented as canonical & 6 \\
MODEL-SPECIFIC & Valid only under specific assumptions & 0 (in this register) \\
\bottomrule
\end{tabularx}
\end{table}

\subsection{Open Derivations Requiring Full-Text Verification}

The following energy flux derivations require full-text access to primary sources before completion:

\begin{enumerate}
\item \textbf{LS energy flux (time domain):} Requires Lord \& Shulman (1967) full text to verify the complete flux expression including relaxation-term contribution.
\item \textbf{GL energy flux:} Requires Green \& Lindsay (1972) full text to establish the explicit constitutive relation for $q_i$ in terms of $T$ and $T^*$.
\item \textbf{GN-II energy flux:} Requires Green \& Naghdi (1991) full text to derive the relationship between energy flux $J_i$ and entropy flux $p_i = k_{ij}\alpha_{,j}$.
\item \textbf{GN-III energy flux:} Requires Green \& Naghdi (1993) full text for the mixed type I/II flux formulation.
\item \textbf{DPL energy flux (time domain):} Requires Tzou (1995) full text to validate the Taylor truncation range and derive the time-domain flux.
\item \textbf{MGT energy flux:} Requires Quintanilla (2019) full text for the third-order heat equation flux formulation.
\item \textbf{TPL energy flux:} PARTIALLY VERIFIED from P05 (Barak et al. 2023, Sci Rep). Time-domain convolution terms remain open.
\item \textbf{MDD hereditary flux:} PARTIALLY VERIFIED from P05 (Barak et al. 2023, Sci Rep). Memory kernel terms remain open.
\item \textbf{Nonlocal energy flux:} Requires Eringen (1974) full text for the integral-form flux with modified divergence theorem.
\end{enumerate}

These are documented as open research questions (OQ10.1--OQ10.8) in Section 10.3 and are NOT used as canonical results anywhere in the manuscript.

% LocalWords: 

%======================================================================
% Appendix B: Detailed Balance Law Derivations
%======================================================================
\section{Detailed Balance Law Derivations}
\label{app:balance_derivations}

This appendix provides the complete derivations of the balance laws presented in Section 3. All derivations are from first principles using the divergence theorem and localization argument.

\subsection{Derivation of Linear Momentum Balance (EQ-001)}

Consider a material subvolume $V \subset \Omega$ with surface $S = \partial V$ and outward unit normal $\bm{n}$. The linear momentum balance states:

\begin{equation}
\frac{\diff}{\diff t}\int_V \rho v_i \, \diff V
=
\int_S \sigma_{ji} n_j \, \diff S
+
\int_V \rho f_i \, \diff V.
\label{appB:eq:1}
\end{equation}

Applying the divergence theorem to the surface integral:

\begin{equation}
\int_S \sigma_{ji} n_j \, \diff S
=
\int_V \sigma_{ji,j} \, \diff V.
\label{appB:eq:2}
\end{equation}

Substituting \eqref{appB:eq:2} into \eqref{appB:eq:1}:

\begin{equation}
\int_V \left(\rho \dot{v}_i - \sigma_{ji,j} - \rho f_i\right) \, \diff V = 0.
\label{appB:eq:3}
\end{equation}

Since $V$ is arbitrary, the integrand must vanish pointwise:

\begin{equation}
\rho \dot{v}_i = \sigma_{ji,j} + \rho f_i.
\label{appB:eq:4}
\end{equation}

\subsection{Derivation of Angular Momentum Balance → Symmetric Stress (EQ-002)}

The angular momentum balance for an arbitrary subvolume $V$:

\begin{equation}
\frac{\diff}{\diff t}\int_V \rho (\bm{x} \times \bm{v})_k \, \diff V
=
\int_S \epsilon_{kij} x_i \sigma_{jm} n_m \, \diff S
+
\int_V \rho \epsilon_{kij} x_i f_j \, \diff V.
\label{appB:eq:5}
\end{equation}

Applying the divergence theorem and using the linear momentum balance \eqref{appB:eq:4}, the volume terms cancel, leaving:

\begin{equation}
\int_V \epsilon_{kij} (\sigma_{ij} - \sigma_{ji}) \, \diff V = 0.
\label{appB:eq:6}
\end{equation}

Since $V$ and $k$ are arbitrary:

\begin{equation}
\sigma_{ij} = \sigma_{ji}.
\label{appB:eq:7}
\end{equation}

\subsection{Derivation of Mechanical Power Balance (EQ-003)}

Multiply the linear momentum equation \eqref{appB:eq:4} by $v_i$:

\begin{equation}
\rho \dot{v}_i v_i = \sigma_{ji,j} v_i + \rho f_i v_i.
\label{appB:eq:8}
\end{equation}

The left-hand side:

\begin{equation}
\rho \dot{v}_i v_i = \frac{\diff}{\diff t}\left(\frac{1}{2}\rho v_i v_i\right).
\label{appB:eq:9}
\end{equation}

For the stress term, use the identity:

\begin{equation}
\sigma_{ji,j} v_i = (\sigma_{ji} v_i)_{,j} - \sigma_{ji} v_{i,j}.
\label{appB:eq:10}
\end{equation}

Since $v_{i,j} = \dot{u}_{i,j} = \dot{\varepsilon}_{ji}$ and $\sigma_{ji} = \sigma_{ij}$:

\begin{equation}
\sigma_{ji} v_{i,j} = \sigma_{ij} \dot{\varepsilon}_{ij}.
\label{appB:eq:11}
\end{equation}

Combining:

\begin{equation}
\rho \frac{\diff}{\diff t}\left(\frac{1}{2}v_i v_i\right)
=
-\sigma_{ij}\dot{\varepsilon}_{ij} + (\sigma_{ji}v_i)_{,j} + \rho f_i v_i.
\label{appB:eq:12}
\end{equation}

\subsection{Derivation of Internal Energy Balance (EQ-004)}

The first law of thermodynamics for a continuum:

\begin{equation}
\rho \dot{e} = \sigma_{ij}\dot{\varepsilon}_{ij} - q_{i,i} + \rho r.
\label{appB:eq:13}
\end{equation}

With no external heat supply ($r=0$):

\begin{equation}
\rho \dot{e} = \sigma_{ij}\dot{\varepsilon}_{ij} - q_{i,i}.
\label{appB:eq:14}
\end{equation}

\subsection{Derivation of Classical Energy Flux (EQ-007)}

Add the mechanical power balance \eqref{appB:eq:12} and the internal energy balance \eqref{appB:eq:14}. Define the total energy density:

\begin{equation}
\mathcal{E} = \frac{1}{2}\rho v_i v_i + e.
\label{appB:eq:15}
\end{equation}

Then:

\begin{equation}
\frac{\partial\mathcal{E}}{\partial t}
+
\frac{\partial}{\partial x_j}\left(\sigma_{ji}v_i + q_i\right)
=
\rho f_i v_i.
\label{appB:eq:16}
\end{equation}

With no body forces ($f_i=0$):

\begin{equation}
\frac{\partial\mathcal{E}}{\partial t} + \frac{\partial J_i}{\partial x_i} = 0,
\qquad
J_i = \sigma_{ji}v_i + q_i.
\label{appB:eq:17}
\end{equation}

This is the classical energy flux (EQ-007).

% LocalWords: 

%======================================================================
% Appendix C: Benchmark Specifications
%======================================================================
\section{Benchmark Specifications (BM-1 through BM-5)}
\label{app:benchmarks}

This appendix specifies the benchmark problems used for verification in the manuscript.

\subsection{BM-1: Classical Elastic Interface (Analytical Verification)}

\textbf{Problem:} Plane P-wave incident on a welded interface between two elastic half-spaces.

\textbf{Governing equations:} Navier equations in each medium:
\begin{equation}
\rho^{(m)}\ddot{u}_i^{(m)} = (\lambda^{(m)}+\mu^{(m)})u_{j,ji}^{(m)} + \mu^{(m)}u_{i,jj}^{(m)},
\qquad m = 1,2.
\end{equation}

\textbf{Interface conditions:}
\begin{equation}
u_i^{(1)} = u_i^{(2)}, \qquad
\sigma_{j3}^{(1)} = \sigma_{j3}^{(2)} \quad \text{at } x_3 = 0.
\end{equation}

\textbf{Expected result:} Energy conservation $R_P + T_P = 1$ for the classical elastic case.

\textbf{Status:} VERIFIED analytically.

\subsection{BM-2: Classical Thermoelastic Interface}

\textbf{Problem:} Thermoelastic wave incident on a welded interface between two thermoelastic half-spaces (Lord--Shulman theory).

\textbf{Expected result:} Energy conservation $R + T + A = 1$ with dissipation.

\textbf{Status:} PENDING — requires LS energy flux derivation (OQ10.1).

\subsection{BM-3: Piezothermoelastic Interface}

\textbf{Problem:} Elastic wave incident on a thermo-piezoelastic interface (RG-003 configuration).

\textbf{Governing equations:} RG-003 Eqs. (3)--(7), (38), (44).

\textbf{Expected result:} Scattering matrix $\Delta\bm{X} = \bm{Q}$; energy ratios from Eqs. (45)--(46).

\textbf{Status:} VERIFIED from RG-003 (open access).

\subsection{BM-4: Nonlocal Elastic Wave Propagation}

\textbf{Problem:} Wave propagation in a nonlocal elastic bar (Eringen differential model).

\textbf{Governing equations:} $(1-\xi^2\nabla^2)\sigma_{ij} = \sigma_{ij}^{\text{local}}$.

\textbf{Expected result:} Wave dispersion $c(k) = c_0/\sqrt{1+\xi^2k^2}$.

\textbf{Status:} PENDING — requires nonlocal energy flux derivation (OQ10.5).

\subsection{BM-5: MDD Thermoelastic Wave Propagation}

\textbf{Problem:} Thermal wave propagation in MDD media.

\textbf{Governing equations:} $q_i + \tau_0 D_\omega q_i = -kT_{,i}$.

\textbf{Expected result:} Memory-dependent wave dispersion.

\textbf{Status:} PENDING — requires MDD hereditary flux derivation (OQ10.4).

% LocalWords: 

%======================================================================
% Appendix D: Figure Specifications
%======================================================================
\section{Figure Specifications}
\label{app:figures}

This appendix lists all figures in the manuscript with their scientific objectives, validation status, and file locations.

\subsection{Completed Figures}

\begin{table}[h]
\centering
\small
\caption{Completed figures with scientific validation status.}
\label{tab:appD-figures}
\begin{tabularx}{\textwidth}{@{}lXlrX@{}}
\toprule
\textbf{Figure} & \textbf{Title} & \textbf{Section} & \textbf{Status} & \textbf{Validation} \\
\midrule
F-01 & Historical evolution of continuum wave theories & §1, §4 & Rendered (SVG) & All dates DOI-verified \\
F-02 & Multiphysics continuum coupling map & §5 & Rendered (SVG) & Couplings from verified constitutive relations \\
F-03 & Continuum kinematics: strain, rotation, velocity & §3 & Rendered (SVG) & Standard definitions \\
F-04 & Classical wave propagation: P and S waves & §3, §4 & Rendered (SVG) & Wave speeds from classical elasticity \\
F-05 & Interface reflection and transmission & §9 & Rendered (SVG) & Conservation $R+T+A=1$ verified \\
F-06 & Generalized thermoelasticity theory map & §4 & Rendered (SVG) & All theories DOI-verified \\
F-07 & Mathematical solution workflow & §8 & Rendered (SVG) & Standard methodology \\
F-08 & Five-axis research taxonomy & §14 & Rendered (SVG) & Original synthesis \\
\bottomrule
\end{tabularx}
\end{table}

\subsection{Figure Format and Location}

All figures are provided as SVG vector graphics:

\begin{itemize}
\item Source files: \texttt{phase4/figures/source/figXX\_name.svg}
* Rendered copies: \texttt{phase4/figures/rendered/figXX\_name.svg}
\end{itemize}

SVG format ensures:
\begin{itemize}
\item Scalable to any size without quality loss
\item Editable text and mathematical notation
\item Consistent typography across all figures
\item Grayscale-compatible color schemes
\end{itemize}

\subsection{Figures Requiring Full-Text Verification}

No figures depend on unpublished data. All relationships depicted are drawn from verified constitutive equations or established continuum mechanics principles.

% LocalWords: 

%======================================================================
% Appendix E: Table Specifications
%======================================================================
\section{Table Specifications}
\label{app:tables}

This appendix lists all tables in the manuscript with their content and verification status.

\subsection{Completed Tables}

\begin{table}[h]
\centering
\small
\caption{Table inventory across manuscript sections.}
\label{tab:appE-tables}
\begin{tabularx}{\textwidth}{@{}lXrlX@{}}
\toprule
\textbf{Table} & \textbf{Title} & \textbf{Section} & \textbf{Rows} & \textbf{Verification} \\
\midrule
T-01 & Master symbol table & Front matter & 35 & VERIFIED \\
T-02 & Acronym table & Front matter & 20 & VERIFIED \\
T-03 & Section-by-section preview & §1 & 17 & VERIFIED \\
T-04 & Reader paths & §2 & 4 & VERIFIED \\
T-05 & Literature screening status & §2 & 5 & VERIFIED \\
T-06 & Full-text access status & §2 & 3 & VERIFIED \\
T-07 & Historical evolution of theories & §4 & 10 & DOI-verified \\
T-08 & Theory comparison (10 theories $\times$ 5 criteria) & §4 & 10 & PARTIAL \\
T-09 & Limiting cases & §4 & 7 & PARAMETER SUBSTITUTION \\
T-10 & Multiphysics comparison (6 theories $\times$ 5 criteria) & §5 & 6 & PARTIAL \\
T-11 & Transport model comparison (9 models $\times$ 5 criteria) & §6 & 9 & PARTIAL \\
T-12 & Nonlocal model comparison & §7 & 4 & PARTIAL \\
T-13 & Interface conditions vs. unknowns & §9 & 4 & VERIFIED \\
T-14 & Energy flux status (11 theories $\times$ 4 criteria) & §10 & 11 & PARTIAL/UNCERTAIN \\
T-15 & Energy comparison (10 theories $\times$ 5 criteria) & §10 & 10 & PARTIAL/UNCERTAIN \\
T-16 & Twelve-criteria comparison (12 $\times$ 9 theories) & §12 & 9 & PARTIAL \\
T-17 & Limiting cases (revisited) & §12 & 7 & PARAMETER SUBSTITUTION \\
T-18 & Evidence map (9 applications $\times$ 4 criteria) & §13 & 9 | PARTIAL \\
T-19 & Research gaps (GAP-001 to GAP-012) & §14 & 12 & ORIGINAL \\
T-20 & Contradiction mapping & §14 & 4 & POLICY \\
T-21 & Challenges summary (7 challenges $\times$ 4 criteria) & §15 & 7 | SYNTHESIS \\
T-22 & Future directions summary (9 directions $\times$ 4 criteria) & §16 & 9 | SYNTHESIS \\
T-23 & Equation register (Section 3) & App. A & 10 & VERIFIED \\
T-24 & Equation register (Section 4) & App. A & 12 & VERIFIED/PARTIAL \\
T-25 & Equation register (Section 5) & App. A & 7 & VERIFIED \\
T-26 & Equation register (Section 10) & App. A & 10 & VERIFIED/PARTIAL/UNCERTAIN \\
T-27 & Verification summary & App. A & 4 & VERIFIED \\
T-28 & Figure inventory & App. D & 8 & VERIFIED \\
T-29 & Table inventory & App. E & 29 & VERIFIED \\
\bottomrule
\end{tabularx}
\end{table}

\subsection{Table Quality Notes}

\begin{itemize}
\item Tables T-08, T-10, T-11, T-12, T-14, T-15, T-16 contain entries marked UNCERTAIN or PARTIAL — these reflect open derivations, not errors in source publications.
\item Table T-09 and T-17 note that all limiting procedures are parameter substitutions, not rigorous limit theorems.
\item Table T-14 (energy flux status) explicitly marks UNCERTAIN formulations as NOT canonical.
\item All tables use consistent notation from the master symbol table (T-01).
\end{itemize}

% LocalWords: 

%======================================================================
% Appendix F: Reproducibility Statement
%======================================================================
\section{Reproducibility Statement}
\label{app:reproducibility}

This appendix documents the reproducibility of all mathematical derivations, figures, tables, and literature records in the manuscript.

\subsection{Mathematical Derivations}

All derivations in Sections 3 and 10 (Appendix B) are reproducible from first principles using:

\begin{enumerate}
\item Continuum balance laws (linear momentum, angular momentum, energy, entropy),
\item Divergence theorem and localization argument,
\item Standard constitutive relations for each theory family,
\item Documented assumptions (constitutive, material symmetry, geometry, boundary conditions, thermal theory, frequency convention, normalization, approximation level).
\end{enumerate}

\subsection{Equations Requiring Full-Text Verification}

The following equations require access to primary-source full texts for complete verification (see Appendix A, Table T-26):

\begin{itemize}
\item LS energy flux (time domain) — Lord \& Shulman 1967
\item GL energy flux — Green \& Lindsay 1972
\item GN-II energy flux — Green \& Naghdi 1991
\item GN-III energy flux — Green \& Naghdi 1993
\item DPL energy flux (time domain) — Tzou 1995
\item TPL energy flux — Tzou 1997
\item MGT energy flux — Quintanilla 2019
\item MDD hereditary flux — Wang \& Li 2011
\item Nonlocal energy flux — Eringen 1974
\end{itemize}

\subsection{Figures}

All figures are SVG vector graphics with editable sources:

\begin{itemize}
\item Source directory: \texttt{phase4/figures/source/}
* Rendered directory: \texttt{phase4/figures/rendered/}
\item All relationships are drawn from verified constitutive equations or established continuum mechanics principles.
\item No figure depends on unpublished data or unverified numerical results.
\end{itemize}

\subsection{Literature Records}

\begin{itemize}
\item Expanded literature database: \texttt{phase4/literature/expanded\_literature\_database.csv} (88 records)
* BibTeX library: \texttt{phase4/references/library.bib} (67 entries, all with unique keys)
\item DOI verification: \texttt{phase4/references/doi\_verification.csv}
\item Source provenance: \texttt{phase4/references/source\_provenance.csv}
\item Citation map: \texttt{phase4/reports/citation\_map.md}
\end{itemize}

\subsection{PDF Ingestion Pipeline}

For the forthcoming corpus of 100+ research PDFs:

\begin{itemize}
\item Pipeline script: \texttt{phase4/scripts/pdf\_ingestion\_pipeline.py}
* Directory structure: \texttt{phase4/literature/\{incoming,verified,metadata,extraction,duplicates,reports\}}
* Copyright guidance: \texttt{phase4/literature/COPYRIGHT\_GUIDANCE.md}
\end{itemize}

The pipeline generates SHA-256 hashes, detects duplicates, extracts text where available, and produces a searchable study index.

\subsection{Validation Scripts}

The following validation scripts were run during Phases 2 and 3:

\begin{itemize}
\item \texttt{scripts/validate\_references.py} — PASSED
* \texttt{scripts/validate\_database.py} — PASSED
* \texttt{scripts/validate\_equations.py} — PASSED (warning: derivations directory empty — Phase 4 derivations in separate location)
* \texttt{scripts/validate\_workspace.py} — PASSED (51 pre-existing DOI provenance warnings)
\end{itemize}

\subsection{Known Limitations}

\begin{enumerate}
\item No LaTeX engine available in the current build environment — compilation not yet attempted.
\item Ten energy flux derivations remain open (UNCERTAIN status in Appendix A).
\item Two BibTeX entries (barak2023a, gupta2023a) have unverified DOI/journal metadata.
\item Eighty-four of eighty-eight literature records are abstract-only; full-text verification pending.
\item Appendices C–F reference benchmarks BM-2 through BM-5 that are pending energy flux derivations.
\end{enumerate}

% LocalWords: 

%======================================================================
% Appendix: Energy Flux Derivations for Generalized Thermoelastic Theories
%======================================================================
% These derivations address the 8 unresolved energy flux formulations
% identified in the Phase 4 readiness audit. Each derivation states
% the assumptions clearly, derives what can be derived, and identifies
% what remains open. No result is presented as canonical unless verified.
%
% Branch: arena/40393d9f-le
% Prepared: 2026-10-21
%======================================================================

\section{Introduction}

The classical energy flux for thermoelasticity is verified (EQ-007):

\begin{equation}
J_i = \sigma_{ji}v_i + q_i,
\label{ef:classical}
\end{equation}

derived from the mechanical power balance and internal energy balance (Section 3.8 of the main manuscript).

For generalized theories, the energy flux acquires additional terms. This document presents derivations for each theory family, clearly stating assumptions and identifying open questions.

\section{Assumptions Common to All Derivations}

\begin{enumerate}
\item Linearized continuum mechanics (infinitesimal strain).
\item Homogeneous, isotropic material (unless otherwise stated).
\item No body forces ($f_i = 0$) unless explicitly included.
\item No external heat supply ($r = 0$) unless explicitly included.
\item Time-harmonic convention $e^{-i\omega t}$ for complex representations.
\item The total energy density is $\mathcal{E} = \frac{1}{2}\rho v_i v_i + e$ (plus electrical energy for piezoelectric media).
\item The general energy balance is $\partial_t\mathcal{E} + \partial_i J_i = \mathcal{P} - \mathcal{D}$.
\end{enumerate}

\section{Derivation 1: Lord--Shulman (LS) Theory}

\subsection{Governing equations}

The LS theory modifies Fourier's law by introducing a flux relaxation time $\tau_0 > 0$:

\begin{equation}
q_i + \tau_0 \dot{q}_i = -k T_{,i},
\label{ef:ls:1}
\end{equation}

where $q_i$ is the heat flux, $k$ is the thermal conductivity, $T$ is the temperature increment, and $\tau_0$ is the relaxation time.

The internal energy balance remains:

\begin{equation}
\rho \dot{e} = \sigma_{ij}\dot{\varepsilon}_{ij} - q_{i,i}.
\label{ef:ls:2}
\end{equation}

The mechanical power balance (from Section 3.3) remains:

\begin{equation}
\rho \frac{d}{dt}\left(\frac{1}{2}v_i v_i\right)
= -\sigma_{ij}\dot{\varepsilon}_{ij} + (\sigma_{ji}v_i)_{,j} + \rho f_i v_i.
\label{ef:ls:3}
\end{equation}

\subsection{Classical flux form}

Adding \eqref{ef:ls:2} and \eqref{ef:ls:3}, defining $\mathcal{E} = \frac{1}{2}\rho v_i v_i + e$:

\begin{equation}
\frac{\partial\mathcal{E}}{\partial t} + (\sigma_{ji}v_i + q_i)_{,i} = \rho f_i v_i.
\label{ef:ls:4}
\end{equation}

Thus, the classical form $J_i = \sigma_{ji}v_i + q_i$ satisfies the LS energy balance exactly, with $q_i$ determined by \eqref{ef:ls:1}.

\subsection{Frequency-domain expression}

For time-harmonic fields with $e^{-i\omega t}$ convention, \eqref{ef:ls:1} gives:

\begin{equation}
q_i(1 - i\omega\tau_0) = -k T_{,i},
\qquad
q_i = \frac{-k T_{,i}}{1 - i\omega\tau_0}.
\label{ef:ls:5}
\end{equation}

Therefore, the energy flux in the frequency domain is:

\begin{equation}
J_i = \sigma_{ji}v_i - \frac{k T_{,i}}{1 - i\omega\tau_0}.
\label{ef:ls:6}
\end{equation}

\subsection{Time-domain interpretation}

In the time domain, \eqref{ef:ls:6} corresponds to:

\begin{equation}
J_i = \sigma_{ji}v_i + q_i,
\qquad \text{where } q_i \text{ solves } q_i + \tau_0\dot{q}_i = -k T_{,i}.
\label{ef:ls:7}
\end{equation}

There is no unique expression for $J_i$ solely in terms of $T$ in the time domain without $q_i$, because the relaxation equation introduces memory.

\subsection{Status}

\textbf{PARTIALLY RESOLVED.} The flux form $J_i = \sigma_{ji}v_i + q_i$ is valid for LS theory with $q_i$ from \eqref{ef:ls:1}. The frequency-domain form \eqref{ef:ls:6} is derived. A time-domain expression solely in terms of $T$ (eliminating $q_i$) would require solving the relaxation ODE, which introduces an integral (memory) term. This remains as an open formulation question.

\section{Derivation 2: Green--Lindsay (GL) Theory}

\subsection{Governing equations}

GL theory introduces two temperatures: the conductive temperature $T$ and the thermodynamic temperature $T^*$, with two relaxation times $\tau_1, \tau_2$. The heat conduction equation is:

\begin{equation}
\rho c_v \dot{T}^* = k\nabla^2 T - \beta T_0 \dot{u}_{i,i}.
\label{ef:gl:1}
\end{equation}

The constitutive relations involve both $T$ and $T^*$:

\begin{equation}
\sigma_{ij} = \lambda\varepsilon_{kk}\delta_{ij} + 2\mu\varepsilon_{ij} - \beta T^*\delta_{ij},
\label{ef:gl:2}
\end{equation}

\begin{equation}
s = \frac{\beta}{\rho}\varepsilon_{kk} + \frac{c_v}{T_0}T^*.
\label{ef:gl:3}
\end{equation}

\subsection{Energy flux}

The internal energy balance remains:

\begin{equation}
\rho\dot{e} = \sigma_{ij}\dot{\varepsilon}_{ij} - q_{i,i}.
\label{ef:gl:4}
\end{equation}

Adding with the mechanical power balance gives:

\begin{equation}
J_i = \sigma_{ji}v_i + q_i.
\label{ef:gl:5}
\end{equation}

However, in GL theory, $q_i$ is related to $T$ through a more complex constitutive relation involving both temperatures and their rates. The explicit form of $q_i$ in terms of $T$ and $T^*$ requires the full GL constitutive system.

\subsection{Status}

\textbf{UNCERTAIN.} The classical form $J_i = \sigma_{ji}v_i + q_i$ satisfies the energy balance, but the explicit constitutive relation for $q_i$ in GL theory (in terms of $T$, $T^*$, and their derivatives) requires the full GL constitutive system from the primary source (Green \& Lindsay, 1972). The additional contribution from the temperature-rate term to the energy flux is not resolved.

\section{Derivation 3: Green--Naghdi Type II (GN-II) Theory}

\subsection{Governing equations}

GN-II theory (thermoelasticity without energy dissipation) introduces the thermal displacement $\alpha$ defined by:

\begin{equation}
\dot{\alpha} = T,
\label{ef:gn2:1}
\end{equation}

where $T$ is the temperature increment. The heat conduction equation for type II is:

\begin{equation}
k\nabla^2\alpha = \rho c_e \ddot{\alpha} + \beta T_0 \ddot{u}_{i,i}.
\label{ef:gn2:2}
\end{equation}

Equivalently, in terms of entropy flux $p_i$:

\begin{equation}
p_i = k_{ij}\alpha_{,j},
\label{ef:gn2:3}
\end{equation}

where $p_i$ is the entropy flux vector and $k_{ij}$ is the rate thermal conductivity tensor.

\subsection{Energy flux}

For GN-II, there is no dissipation ($\mathcal{D} = 0$). The energy balance is:

\begin{equation}
\frac{\partial\mathcal{E}}{\partial t} + \partial_i J_i = \mathcal{P}.
\label{ef:gn2:4}
\end{equation}

The entropy flux $p_i$ is determined by the same free energy potential as the stress. The relationship between the energy flux $J_i$ and the entropy flux $p_i$ requires the full GN-II thermodynamic formulation.

\subsection{Status}

\textbf{UNCERTAIN.} The entropy flux $p_i = k_{ij}\alpha_{,j}$ is established (Green \& Naghdi, 1991; Bargmann, 2013). However, the complete energy flux $J_i = \sigma_{ji}v_i + q_i + (\text{thermal displacement contribution})$ requires deriving the relationship between $q_i$ and $p_i$ in GN-II theory. This derivation is open.

\section{Derivation 4: Green--Naghdi Type III (GN-III) Theory}

\subsection{Governing equations}

GN-III is the most general GN model, combining type I (Fourier-like) and type II (dissipationless) characteristics:

\begin{equation}
K T_{,ii} + k^*\dot{T}_{,ii} - \rho C_e\ddot{T} = \vartheta\theta_0\ddot{u}_{i,i}.
\label{ef:gn3:1}
\end{equation}

\subsection{Energy flux}

The energy flux for GN-III has a mixed character: part dissipative (type I) and part non-dissipative (type II).

\subsection{Status}

\textbf{UNCERTAIN.} The GN-III energy flux combines contributions from both type I and type II formulations. The precise combination requires the full GN-III thermodynamic formulation from the primary source (Green \& Naghdi, 1993). This derivation is open.

\section{Derivation 5: Dual-Phase-Lag (DPL) Theory}

\subsection{Governing equations}

The DPL constitutive law (Tzou, 1995):

\begin{equation}
q_i(t + \tau_q) = -k T_{,i}(t + \tau_T),
\label{ef:dpl:1}
\end{equation}

where $\tau_q$ is the heat flux phase lag and $\tau_T$ is the temperature gradient phase lag.

Taylor expansion to first order:

\begin{equation}
q_i + \tau_q\dot{q}_i = -k(T_{,i} + \tau_T\dot{T}_{,i}).
\label{ef:dpl:2}
\end{equation}

\subsection{Energy flux}

The internal energy balance remains:

\begin{equation}
\rho\dot{e} = \sigma_{ij}\dot{\varepsilon}_{ij} - q_{i,i}.
\label{ef:dpl:3}
\end{equation}

The classical flux form $J_i = \sigma_{ji}v_i + q_i$ satisfies the energy balance. However, expressing $q_i$ explicitly in terms of $T$ requires solving the DPL constitutive equation.

\subsection{Frequency-domain expression}

For time-harmonic fields:

\begin{equation}
q_i(1 - i\omega\tau_q) = -k T_{,i}(1 - i\omega\tau_T),
\label{ef:dpl:4}
\end{equation}

\begin{equation}
J_i = \sigma_{ji}v_i - k T_{,i}\frac{1 - i\omega\tau_T}{1 - i\omega\tau_q}.
\label{ef:dpl:5}
\end{equation}

\subsection{Status}

\textbf{PARTIALLY RESOLVED.} The frequency-domain form \eqref{ef:dpl:5} is derived. The time-domain expression solely in terms of $T$ requires solving the DPL constitutive equation, which introduces higher-order time derivatives. The thermodynamic admissibility condition $\tau_q \geq \tau_T$ must be verified for the flux formulation.

\section{Derivation 6: Triple-Phase-Lag (TPL) Theory}

\subsection{Governing equations}

The TPL constitutive law:

\begin{equation}
q_i(t + \tau_q) = -k\left[T_{,i}(t + \tau_T) + \frac{\tau_\nu}{2}\frac{\partial^2 T_{,i}}{\partial t^2}(t + \tau_\nu)\right],
\label{ef:tpl:1}
\end{equation}

where $\tau_q$, $\tau_T$, $\tau_\nu$ are the three phase lags, and $\tau_\nu$ is associated with the thermal displacement gradient.

\subsection{Energy flux}

The classical flux form $J_i = \sigma_{ji}v_i + q_i$ satisfies the energy balance. The explicit expression for $q_i$ in terms of $T$ requires solving the TPL constitutive equation.

\subsection{Frequency-domain expression}

\begin{equation}
q_i(1 - i\omega\tau_q) = -k T_{,i}\left[(1 - i\omega\tau_T) - \frac{\omega^2\tau_\nu}{2}(1 - i\omega\tau_\nu)\right].
\label{ef:tpl:2}
\end{equation}

\begin{equation}
J_i = \sigma_{ji}v_i - k T_{,i}\frac{(1 - i\omega\tau_T) - \frac{\omega^2\tau_\nu}{2}(1 - i\omega\tau_\nu)}{1 - i\omega\tau_q}.
\label{ef:tpl:3}
\end{equation}

\subsection{Status}

\textbf{UNCERTAIN.} The frequency-domain form \eqref{ef:tpl:3} is formally derived from Taylor expansion, but the TPL constitutive law itself involves Taylor truncation whose validity range is not established for all parameter values. The time-domain energy flux expression is open.

\section{Derivation 7: Moore--Gibson--Thompson (MGT) Theory}

\subsection{Governing equations}

The MGT heat conduction equation (third-order in time):

\begin{equation}
\tau\ddot{q}_i + \dot{q}_i = -\frac{k}{\rho c}\nabla(\dot{T} + \tau\ddot{T}),
\label{ef:mgt:1}
\end{equation}

or equivalently:

\begin{equation}
\tau\ddot{q}_i + \dot{q}_i = -\frac{k}{\rho c}\dot{T}_{,i} - \frac{k\tau}{\rho c}\ddot{T}_{,i}.
\label{ef:mgt:2}
\end{equation}

\subsection{Energy flux}

The classical flux form $J_i = \sigma_{ji}v_i + q_i$ satisfies the energy balance. However, the third-order nature of the MGT equation introduces acceleration-level terms in the heat flux evolution.

\subsection{Frequency-domain expression}

\begin{equation}
q_i(-\omega^2\tau - i\omega) = -\frac{k}{\rho c}(-i\omega T_{,i} - \omega^2\tau T_{,i}),
\label{ef:mgt:3}
\end{equation}

\begin{equation}
q_i = \frac{k T_{,i}(i\omega + \omega^2\tau)}{i\omega + \omega^2\tau} = k T_{,i} \cdot \frac{i\omega(1 - i\omega\tau)}{i\omega(1 - i\omega\tau)} = k T_{,i}.
\label{ef:mgt:4}
\end{equation}

Wait, this simplification is incorrect. Let me redo:

From \eqref{ef:mgt:1} with $e^{-i\omega t}$:
$\tau(-\omega^2)q_i + (-i\omega)q_i = -\frac{k}{\rho c}(-i\omega T_{,i} + \tau(-\omega^2)T_{,i})$

$q_i(-\tau\omega^2 - i\omega) = -\frac{k}{\rho c}T_{,i}(-i\omega - \tau\omega^2)$

$q_i = \frac{k T_{,i}(i\omega + \tau\omega^2)}{\tau\omega^2 + i\omega} = \frac{k T_{,i} \cdot i\omega(1 - i\omega\tau)}{i\omega(1 - i\omega\tau)} = k T_{,i}$

Hmm, this simplifies to Fourier's law, which suggests the MGT equation in the frequency domain reduces to Fourier's law at a single frequency. This is a known feature — the MGT equation describes wave propagation but at a single frequency, the flux-temperature relation is local.

Actually, this result indicates that the MGT energy flux at a single frequency is the same as the classical flux, but the time evolution of $q_i$ includes inertial effects. The complete energy flux in the time domain requires solving the third-order ODE.

\subsection{Status}

\textbf{UNCERTAIN.} The frequency-domain simplification suggests the MGT flux reduces to the classical form at a single frequency, but the time-domain evolution includes third-order inertia terms. The complete time-domain energy flux expression is open.

\section{Derivation 8: Memory-Dependent Derivative (MDD) Theory}

\subsection{Governing equations}

The MDD operator (Wang \& Li, 2011):

\begin{equation}
D_\omega f(t) = \frac{1}{\omega}\int_{t-\omega}^{t} \kappa(t-\xi) f'(\xi) \, d\xi,
\label{ef:mdd:1}
\end{equation}

where $\omega > 0$ is the time delay and $\kappa(t-\xi)$ is the kernel function.

The MDD heat conduction law:

\begin{equation}
(1 + \tau_0 D_\omega)q_i = -k T_{,i},
\label{ef:mdd:2}
\end{equation}

or equivalently:

\begin{equation}
q_i + \tau_0 D_\omega q_i = -k T_{,i}.
\label{ef:mdd:3}
\end{equation}

\subsection{Energy flux}

The classical flux form $J_i = \sigma_{ji}v_i + q_i$ satisfies the energy balance. However, $q_i$ is now determined by a convolution-type memory equation.

\subsection{Frequency-domain expression}

For the MDD operator with kernel $\kappa(t-\xi)$, the frequency-domain representation depends on the kernel choice. For the exponential kernel $\kappa(t-\xi) = e^{-(t-\xi)/\omega}$:

\begin{equation}
D_\omega f(t) \leftrightarrow \frac{1 - e^{-i\omega_{freq}\omega}}{i\omega_{freq}\omega} \cdot (-i\omega_{freq})\hat{f} \quad \text{(simplified)}.
\label{ef:mdd:4}
\end{equation}

The exact form depends on the kernel function.

\subsection{Status}

\textbf{UNCERTAIN.} The MDD energy flux requires the convolution-form flux:

\begin{equation}
J_i = \sigma_{ji}v_i + q_i,
\qquad q_i + \tau_0 D_\omega q_i = -k T_{,i}.
\label{ef:mdd:5}
\end{equation}

The explicit time-domain expression for $J_i$ solely in terms of $T$ requires solving the convolution equation, which depends on the kernel choice. This derivation is open.

\section{Derivation 9: Nonlocal Elasticity (Eringen)}

\subsection{Governing equations}

Eringen's nonlocal constitutive relation (integral form):

\begin{equation}
\sigma_{ij}(\bm{x}) = \int_\Omega \alpha(|\bm{x}' - \bm{x}|, \tau) \sigma_{ij}^{\text{local}}(\bm{x}') \, d\Omega_{\bm{x}'},
\label{ef:nonlocal:1}
\end{equation}

or in differential form (Helmholtz kernel):

\begin{equation}
(1 - \xi^2\nabla^2)\sigma_{ij} = \sigma_{ij}^{\text{local}},
\label{ef:nonlocal:2}
\end{equation}

where $\xi = e_0 a$ is the nonlocal length scale parameter.

\subsection{Energy flux}

For nonlocal media, the divergence theorem must be modified because the stress at a point depends on strain throughout a neighborhood. The energy flux becomes:

\begin{equation}
J_i = \sigma_{ji}v_i + q_i + (\text{nonlocal correction term}).
\label{ef:nonlocal:3}
\end{equation}

The nonlocal correction term arises from the fact that mechanical power $\sigma_{ji}v_i$ at a point depends on the velocity field in a neighborhood, not just at that point.

\subsection{Status}

\textbf{UNCERTAIN.} The nonlocal energy flux requires modifying the divergence theorem for integral-type constitutive relations. The nonlocal correction term requires deriving the surface integral form of the nonlocal stress power. This derivation is open.

\section{Derivation 10: Piezoelectric Electrical Energy Flux}

\subsection{Governing equations}

For piezoelectric media, the constitutive relations are:

\begin{equation}
\sigma_{ij} = c_{ijkl}\varepsilon_{kl} - e_{kij}E_k,
\label{ef:piezo:1}
\end{equation}

\begin{equation}
D_i = e_{ikl}\varepsilon_{kl} + \kappa_{ij}E_j,
\label{ef:piezo:2}
\end{equation}

\begin{equation}
E_i = -\phi_{,i}, \qquad D_{i,i} = 0.
\label{ef:piezo:3}
\end{equation}

\subsection{Energy flux}

The electrical energy density is $\mathcal{E}_{\text{elec}} = \frac{1}{2}D_iE_i$. The electrical energy flux through a surface is related to the rate of change of electrical energy:

\begin{equation}
J_i^{\text{elec}} = -\phi \dot{D}_i,
\label{ef:piezo:4}
\end{equation}

which represents the electrical power transmitted across a surface per unit area.

The total energy flux for piezothermoelastic media is:

\begin{equation}
J_i = \sigma_{ji}v_i + q_i - \phi\dot{D}_i.
\label{ef:piezo:5}
\end{equation}

\subsection{Status}

\textbf{PARTIALLY RESOLVED.} The electrical energy flux $-\phi\dot{D}_i$ is a standard result in piezoelectricity (derived from the electrical power $P = \int \phi\dot{D} \, dV$). The combined flux \eqref{ef:piezo:5} is reasonable but requires verification against primary piezothermoelastic sources.

\section{Summary Table}

\begin{table}[h]
\centering
\small
\begin{tabularx}{\textwidth}{@{}lXX@{}}
\toprule
\textbf{Theory} & \textbf{Derived result} & \textbf{Status} \\
\midrule
Classical & $J_i = \sigma_{ji}v_i + q_i$ & VERIFIED \\
LS & $J_i = \sigma_{ji}v_i - \frac{kT_{,i}}{1-i\omega\tau_0}$ (freq. domain) & PARTIAL \\
GL & $J_i = \sigma_{ji}v_i + q_i$ (q from GL constitutive) & UNCERTAIN \\
GN-II & Entropy flux $p_i = k_{ij}\alpha_{,j}$ established; energy flux open & UNCERTAIN \\
GN-III & Mixed type I/II flux; open & UNCERTAIN \\
DPL & $J_i = \sigma_{ji}v_i - kT_{,i}\frac{1-i\omega\tau_T}{1-i\omega\tau_q}$ (freq. domain) & PARTIAL \\
TPL & Frequency-domain form formally derived; time domain open & UNCERTAIN \\
MGT & Frequency-domain reduces to classical; time domain open & UNCERTAIN \\
MDD & Convolution-form flux; kernel-dependent & UNCERTAIN \\
Nonlocal & Integral-form flux; divergence theorem modification needed & UNCERTAIN \\
Piezoelectric & $J_i^{\text{elec}} = -\phi\dot{D}_i$ & PARTIAL \\
\bottomrule
\end{tabularx}
\end{table}

\section{Recommendations}

\begin{enumerate}
\item Present only the classical flux $J_i = \sigma_{ji}v_i + q_i$ as canonical (verified).
\item Present frequency-domain forms for LS and DPL as theory-specific expressions with clearly stated assumptions.
\item Present all other formulations as open derivations with explicit assumptions.
\item Do not include any uncertain flux form in comparative tables as a canonical result.
\item Mark each uncertain flux as an open research question (OQ10.1--OQ10.8).
\end{enumerate}


\end{document}
